\section*{Capítulo \ref{ch:b3:sequence-stratigraphy} --- Estratigrafia de sequências}

\begin{solution}{exo:b3:sequence-stratigraphy:1}
$A = 0.5 + 0.1 = \qty{0.6}{m/kyr}$; em \qty{20}{kyr}, \qty{12}{m} de
acomodação. Com o nível do mar descendo, $A = -0.3 + 0.1 = \qty{-0.2}{m/kyr}$: uma
queda relativa, atenuada pela subsidência.
\end{solution}

\begin{solution}{exo:b3:sequence-stratigraphy:2}
(a) \omterm{def:b3:sequence-stratigraphy:tract}{Trato de sistemas de estágio de queda} (\omterm{def:b3:sequence-stratigraphy:regression}{regressão forçada}). (b) Trato transgressivo. (c)
\omterm{def:b3:sequence-stratigraphy:tract}{Trato de sistemas de mar alto}. (d) \omterm{def:b3:sequence-stratigraphy:tract}{Trato de sistemas de mar baixo} (\omterm{def:b3:sequence-stratigraphy:regression}{regressão normal} após a queda).
\end{solution}

\begin{solution}{exo:b3:sequence-stratigraphy:3}
Cada \omterm{def:b3:sequence-stratigraphy:parasequence}{parassequência} termina mais rasa que a de baixo: um empilhamento progradacional,
com a linha de costa avançando em direção ao mar, como em um \omterm{def:b3:sequence-stratigraphy:tract}{trato de sistemas de mar alto} ou de mar baixo. A
ordem oposta, com cada uma terminando mais fina e mais profunda, é retrogradacional: um
trato transgressivo.
\end{solution}

\begin{solution}{exo:b3:sequence-stratigraphy:4}
Cinco pares por feixe é a razão entre a excentricidade curta (cerca de
\qty{100}{kyr}) e a precessão (cerca de \qty{20}{kyr}): cada par é um
ciclo de precessão, cada feixe é um ciclo de excentricidade. A taxa é
$1.6/100 = \qty{1.6}{cm/kyr}$, e um par de \qty{0.32}{m} dura cerca de
\qty{20}{kyr}, o que é consistente.
\end{solution}

\begin{solution}{exo:b3:sequence-stratigraphy:5}
$\mathrm dx_s/\mathrm dt = (4000 - 0.15\times20\,000)/50 = \qty{20}{m/kyr}$,
em direção ao mar. A linha de costa permanece parada quando $AW = Q$: $A = 4000/20\,000 =
\qty{0.2}{m/kyr}$.
\end{solution}

\begin{solution}{exo:b3:sequence-stratigraphy:6}
$\mathrm dx_s/\mathrm dt = 4000/50 + 0.2/0.01 = 80 + 20 = \qty{100}{m/kyr}$;
a queda, sozinha, responde por \qty{20}{m/kyr}, um quinto.
\end{solution}

\begin{solution}{exo:b3:sequence-stratigraphy:7}
A linha de costa recua quando $Q < AW = 0.4\times25\,000 = \qty{10000}{m^3/kyr}$ por
metro. Com $Q = 6000$: $(6000 - 10\,000)/10 = \qty{-400}{m/kyr}$, um recuo
de \qty{0.4}{km} a cada mil anos: uma antepraia rasa produz uma
transgressão rápida.
\end{solution}

\begin{solution}{exo:b3:sequence-stratigraphy:8}
$1/94.9 - 1/123.9 = \qty{0.00247}{ciclos/kyr}$: um registro mais longo que
$1/0.00247 = \qty{406}{kyr}$ separa os dois, e \qty{600}{kyr} (resolução
0.00167) é suficiente. $1/94.9 - 1/98.9 = 0.00042$: são necessários cerca de
\qty{2400}{kyr}, quatro vezes mais que o registro.
\end{solution}

\begin{solution}{exo:b3:sequence-stratigraphy:9}
$16.4/(20\times40.98) = \qty{0.0200}{m/kyr}$, \qty{2.0}{cm/kyr}; as camadas de cinza
estão separadas por $3.1/0.020 = \qty{155}{kyr}$.
\end{solution}

\begin{solution}{exo:b3:sequence-stratigraphy:10}
(a) $\tan\alpha = -0.4/200 = -0.002$: descendo a \ang{0.11}. (b)
$0.2/50 = 0.004$: subindo a \ang{0.23}. Uma seção sísmica abrange dezenas de
quilômetros com exagero vertical de dezenas a centenas, por isso o ângulo aparece;
um penhasco de algumas centenas de metros de comprimento mostra um desnível de um metro ou menos, perdido na
variação das camadas.
\end{solution}

\begin{solution}{exo:b3:sequence-stratigraphy:11}
$1\,296\,000/54.725282 = \qty{23.68}{kyr}$ ($g_5$); $/57.924383 =
\qty{22.37}{kyr}$ ($g_2$); $/68.377912 = \qty{18.95}{kyr}$ ($g_4$);
$/67.834313 = \qty{19.11}{kyr}$ ($g_3$). $g_2 - g_5 = 3.199101''$ por ano,
\qty{405.1}{kyr}; $g_4 - g_5 = 13.65263''$, \qty{94.9}{kyr}.
\end{solution}

\begin{solution}{exo:b3:sequence-stratigraphy:12}
Ambos produzem $0.3\times100 = \qty{30}{m}$ de acomodação à mesma taxa; com
o mesmo aporte, a geometria é a mesma, já que o sedimento responde apenas ao
nível relativo do mar. Uma mudança eustática é a mesma em toda parte: a mesma
sequência, de mesma idade, em outras margens, aponta para a primeira história. Uma
história de subsidência construída a partir da história térmica e de carga da bacia
(backstripping), ou um registro eustático independente, como isótopos de oxigênio,
também as separa.
\end{solution}

\begin{solution}{pb:b3:sequence-stratigraphy:1}
\textbf{1.} $\dot E = 30\times2\pi/400\times\sin(2\pi t/400)$: no máximo
\qty{0.471}{m/kyr}, subindo em \qty{100}{kyr} e \qty{500}{kyr}, descendo em
\qty{300}{kyr}.

\textbf{2.} \qty{0.020}{m/kyr} na charneira, 0.037 a \qty{20}{km}, 0.062 a
\qty{50}{km}, \qty{0.120}{m/kyr} a \qty{120}{km}.

\textbf{3.} $A < 0$ quando $\sin(2\pi t/400) < -0.0346/0.471 = -0.0734$: $t =
200 + (400/2\pi)\arcsin0.0734 = \qty{204.7}{kyr}$.

\textbf{4.} Com $\sigma = 0.0617$: $t = 400 - 63.66\arcsin0.131 =
\qty{391.6}{kyr}$.

\textbf{5.} Cerca de \qty{187}{kyr}, contra 200: a queda eustática precisa primeiro
superar a subsidência, e deixa de vencê-la antes do mar baixo.

\textbf{6.} $0.045\times400 = \qty{18}{m}$, frente a uma amplitude eustática de
\qty{60}{m}: um terço.

\textbf{7.} Em direção ao continente, de 50 a cerca de \qty{120}{kyr}, em direção ao mar até cerca de
\qty{440}{kyr}, em direção ao continente até cerca de \qty{560}{kyr}, e de novo em direção ao mar até
\qty{600}{kyr} (os pontos de retorno do modelo são 118.5, 439.5 e
\qty{559}{kyr}; uma tabela a cada \qty{50}{kyr} apenas os delimita).

\textbf{8.} $(6000 - 0.30\times15\,000)/40 = \qty{37.5}{m/kyr}$, em direção ao mar.

\textbf{9.} $6000/15\,000 = \qty{0.40}{m/kyr}$.

\textbf{10.} $6000/35.7 + 0.424/0.012 = 168.1 + 35.3 = \qty{203}{m/kyr}$,
dentro de \qty{0.5}{\percent} do modelo.

\textbf{11.} O aporte: \qty{83}{\percent} da velocidade.

\textbf{12.} $(-42.7 - 16.3)/(50.7 - 17.4) = \qty{-1.77}{m/km}$, \ang{0.10}.
Desde então, o primeiro ponto subsidiu $0.0345\times395.5 =
\qty{13.6}{m}$, o segundo $0.0623\times208.5 = \qty{13.0}{m}$: praticamente
o mesmo valor; logo, nas elevações atuais, a descida, \qty{59.0}{m}, difere da
queda eustática apenas em \qty{0.6}{m}. A queda relativa da linha de costa na
época, eustasia menos a subsidência durante a queda, foi menor.

\textbf{13.} 150 e 200: mar alto; 250, 300, 350: estágio de queda; 400:
mar baixo; 450, 500, 550: transgressivo (a linha de costa recua logo após
440); 600: mar alto.

\textbf{14.} 204.7 e \qty{391.6}{kyr} contra 204.5 e 391.5: os mesmos,
dentro do passo do modelo de \qty{0.5}{kyr}.

\textbf{15.} Na plataforma, a superfície exposta durante a queda: o topo dos
depósitos de mar alto, da charneira até \qty{17}{km}, prolongado sobre as
frentes deposicionais em degrau do trato de estágio de queda; na bacia, a
concordância correlata na base ou no topo do trato de estágio de queda, por
convenção. Uma plataforma real seria cortada por rios em \omterm{def:b3:sequence-stratigraphy:condensed}{vales incisos},
intemperizada em solos e parcialmente erodida.

\textbf{16.} $(18.5 - 52.8)/(559 - 439.5) = \qty{-0.29}{km/kyr}$; a
\omterm{def:b3:sequence-stratigraphy:surfaces}{superfície de inundação máxima} se forma ao seu final.

\textbf{17.} A linha de costa varre rapidamente uma planície costeira quase
plana: o aporte fica preso perto dela, espalhado por uma área ampla, e cada
local recebe pouco. Em direção à bacia, além do pé do talude, a cerca de \qty{55}{km}, o
fundo marinho sofre escassez de sedimentos: forma-se ali uma \omterm{def:b3:sequence-stratigraphy:condensed}{seção condensada}.

\textbf{18.} Resolução $1/30 = \qty{0.033}{ciclos/m}$; Nyquist $1/(2\times
0.01) = \qty{50}{ciclos/m}$.

\textbf{19.} $7.5/2.0 = 3.75$, 1 e $0.81/2.0 = 0.405$; as razões orbitais
são $405/100 = 4.05$ e $41/100 = 0.41$.

\textbf{20.} \qty{2.0}{m} para \qty{100}{kyr}: \qty{2.0}{cm/kyr}; $0.81/41 =
\qty{1.98}{cm/kyr}$; $7.5/405 = \qty{1.85}{cm/kyr}$. Os comprimentos de onda
que o \omterm{def:b3:sequence-stratigraphy:periodogram}{periodograma} pode mostrar perto de \qty{8}{m} são 10, 7.5 e \qty{6}{m}:
o valor verdadeiro, \qty{8.1}{m}, cai entre dois deles.

\textbf{21.} $30/0.020 = \qty{1500}{kyr}$.

\textbf{22.} $23.7\times0.020 = \qty{0.474}{m}$. Amostrado a cada
\qty{25}{cm}, a frequência de Nyquist é de 2 ciclos por metro, um comprimento
de onda de \qty{0.5}{m}: a precessão, em 2.11 ciclos por metro, está além dela
e se dobra para $4 - 2.11 = 1.89$, um comprimento de onda falso de \qty{0.53}{m}.

\textbf{23.} Um único pico se ajusta a algum período, em alguma taxa, qualquer
que seja a sua origem; três picos cujas razões coincidem com as razões orbitais
são independentes da taxa, que então deve sair a mesma a partir de cada um.

\textbf{24.} A linha de costa passa de \qty{17.37}{km} no início da queda para
\qty{52.84}{km} no fim da regressão de mar baixo: um deslocamento máximo em
direção à bacia de \textbf{\qty[text-series-to-math=true]{35.5}{km}}.
\end{solution}
